% answers_short.tex —— 2011 年·四川卷·文科 · 简短答案速查
% 由 tools/make_short_answers.py 自 answers.tex 的 \gksolhead/\gkanswer 提取，
% 并与 2011/2011_数学真题_四川文_答案要点.md 逐题比对后生成，勿手改（改源 MD 后重跑）。

\begin{gknotice}
\textbf{使用说明}：本节只给最终答案便于快速核对，完整推导见解析版。答案为编写者独立推导，\textbf{非官方标准答案}。
\end{gknotice}

\gkpart{选择题}
\begin{center}\begin{tabular}{ccccccccccccc}
\toprule
题号 & 1 & 2 & 3 & 4 & 5 & 6 & 7 & 8 & 9 & 10 & 11 & 12 \\
\midrule
答案 & B & B & D & A & A & B & D & C & A & C & A & B \\
\bottomrule
\end{tabular}\end{center}

\gkpart{填空题}
\begin{itemize}[leftmargin=2.6em]
  \item[\textbf{13.}] $84$。
  \item[\textbf{14.}] $16$。
  \item[\textbf{15.}] $32\pi$。
  \item[\textbf{16.}] ②③④。
\end{itemize}

\gkpart{解答题}
\begin{itemize}[leftmargin=2.6em]
  \item[\textbf{17.}] (1) 甲、乙在三小时以上且不超过四小时还车的概率都是 $\dfrac14$；(2) $\dfrac34$
  \item[\textbf{18.}] (1) 最小正周期 $2\pi$，最小值 $-2$；(2) 证明见解析
  \item[\textbf{19.}] (1) 证明见解析；(2) $\dfrac23$
  \item[\textbf{20.}] (1) $q=\dfrac{1+\sqrt5}{2}$ 或 $q=\dfrac{1-\sqrt5}{2}$；(2) 证明见解析
  \item[\textbf{21.}] (1) $\dfrac{16}{7}$；(2) 定值为 $4$
  \item[\textbf{22.}] (1) 单调递增区间为 $[0,2]$，单调递减区间为 $[2,+\infty)$，极大值为 $25$，无极小值；(2) 当 $a\leqslant1$ 或 $a>5$ 时方程无解；当 $1<a\leqslant4$ 时 $x=3-\sqrt{5-a}$；当 $4<a<5$ 时 $x=3\pm\sqrt{5-a}$；当 $a=5$ 时 $x=3$；(3) 证明见解析（$n=1$ 时取等号）
\end{itemize}
